\uuid{1SuC}
\titre{Résolution d'équation différentielle du premier ordre}
\auteur{bodin ; croizat ; sacre}
\organisation{exo7}
\licence{CC BY-NC-SA 4.0}
\datecreate{2022-03-08}
\niveau{L1}
\module{Analyse}
\chapitre{Equation différentielle}
\sousChapitre{Résolution d'équation différentielle du premier ordre}
\difficulte{3}
\theme{}
\exo7qcmid{hilisit.equations-differentielles.055}

\contenu{
\begin{question}
\qclassification{225.01}
Soit $f$ une solution de l'équation différentielle $(H)$ : $y'=4y$. De quelle équation différentielle la fonction $g : x \mapsto f(x)+e^{2x}$ sera-t-elle solution ?
\begin{answers}
	\bad{$y'=4y+e^{2x}$}
	\bad{$y'-4y=4e^{2x}$}
	\good{$y'=4y-2e^{2x}$}
	\bad{$y'=2y$}
\end{answers}
\begin{explanations}
Si $f$ est solution de $(H)$, alors $f'(x)=4f(x)$. On calcule alors la dérivée de $g$ : $g'(x) = f'(x) + 2e^{2x}$. Mais on a alors : $g'(x) = 4f(x) + 2e^{2x} = 4f(x) + 4e^{2x} - 2e^{2x} = 4(f(x)+e^{2x})-2e^{2x} = 4g(x) - 2e^{2x}$. De ce fait, $g$ est solution de $y'=4y-2e^{2x}$. On pouvait aussi exploiter le fait que $f(x)$ s'écrit sous la forme $C e^{4x}$ et tenter de remplacer directement dans chacune des équations différentielles les expressions de $g$ et de $g'$ pour voir quelle égalité était vérifiée.
\end{explanations}
\end{question}
}
